A substituição da especificação técnica 0126 pela especificação~1580 alterou os critérios da Enel Distribuição Ceará para a conexão de edificações com múltiplas unidades consumidoras. Este trabalho compara os requisitos aplicáveis ao Ceará e avalia as consequências de sua aplicação ao projeto elétrico aprovado do Edifício Palatium, em Sobral. A análise documental confrontou as duas especificações com o memorial e as pranchas do empreendimento. Dos 82 requisitos comparados, 39 foram selecionados para o estudo de caso e corresponderam a 38 registros de impacto. Entre esses registros, 23 indicaram adequação potencialmente necessária, sete apontaram compatibilidade, dois dependeram de avaliação da distribuidora e seis permaneceram sem desfecho por insuficiência documental. No cenário estudado, o reposicionamento do centro de proteção geral repercute na localização da medição de incêndio e na proteção contra surtos. A derivação elétrica de incêndio pode ser preservada. A demanda declarada de 291,33~kVA permaneceu como referência, pois os documentos não discriminam a contribuição da recarga veicular para a demanda simultânea nem justificam o multiplicador residencial adotado. Os resultados identificam decisões de projeto que requerem revisão e os dados necessários para resolver os casos pendentes.

\palavraschave{distribuição de energia elétrica; conexão coletiva; demanda elétrica; medição de energia; projeto elétrico.}
